\begin{table}[htbp]
\centering
\caption{The adjusted college gap under inverse-retention-probability weighting}
\label{tab:ipw}
\begin{threeparttable}
\begin{tabular}{lcc}
\toprule
& Survey weights & Survey $\times$ inverse-retention weights \\
Period & (1) & (2) \\
\midrule
Pre-pandemic & -0.009$^{***}$ & -0.009$^{***}$ \\
 & (0.001) & (0.001) \\
Onset & -0.035$^{***}$ & -0.035$^{***}$ \\
 & (0.001) & (0.001) \\
Mid-pandemic & 0.018$^{***}$ & 0.018$^{***}$ \\
 & (0.001) & (0.001) \\
Post-pandemic & -0.010$^{***}$ & -0.010$^{***}$ \\
 & (0.001) & (0.001) \\
\bottomrule
\end{tabular}
\begin{tablenotes}[flushleft]\footnotesize
\item Notes: Both columns are estimated on the matched origins of the retention sample (8,002,128 person-quarters) with the full specification, so the two differ only in the weights. Column (2) multiplies the survey weight by the inverse of a predicted probability of being matched into $t+1$, trimmed at 0.05, from a separate logit per year-quarter on education, demographics, job characteristics and state, occupation and sector fixed effects. Standard errors are two-way clustered by primary sampling unit and year-quarter. Periods are pre-pandemic 2012Q1--2019Q4, onset 2020Q1, mid-pandemic 2020Q2--2021Q3 and post-pandemic 2021Q4--2024Q4.
\end{tablenotes}
\end{threeparttable}
\end{table}
